\documentclass[aspectratio=1610]{beamer}
\usepackage[T1]{fontenc}
\usetheme{wildcat}
\usetikzlibrary{arrows.meta,angles,quotes,calc,intersections,positioning}

\usepackage{amsmath,amssymb,amsfonts}
\usepackage{booktabs}
\usepackage{relsize}
\usepackage{pgfplots}
\pgfplotsset{compat=1.16}
\usepackage{array}
\usepackage{siunitx}
\usepackage{cancel}
\usepackage{jupyter}
\usepackage{minted}
\usepackage{xfrac}
\usepackage{mathtools}

\let\oldfootnotesize\footnotesize
\renewcommand*{\footnotesize}{\oldfootnotesize\tiny}

\def\mathdefault#1{#1}
\everymath=\expandafter{\the\everymath\displaystyle}


\title{Elements of Modern C++ and Fortran \\
       {\small\it ME 701 - Phase III - Lesson 03 - Arrays}}

\date{\input{term.txt} \\ {\footnotesize Git SHA: \input{git_sha.txt}}}

\author{Jeremy Roberts}


\definecolor{ksupurple}{HTML}{512888}
\definecolor{orange}{HTML}{CA7C1B}
\definecolor{skyblue}{RGB}{180,220,255}
\definecolor{light-gray}{gray}{0.97}

\usepackage{listings}

\lstset{basicstyle=\fontsize{7}{8}\selectfont\ttfamily,
        numbers=left,
        numberstyle=\fontsize{5}{6}\selectfont\ttfamily,
        numbersep=5pt,                  
        backgroundcolor=\color{light-gray},
        frame=single, 
        %rulesepcolor=\color{red},
        keywordstyle=\color[rgb]{0,0,1},
        commentstyle=\color[rgb]{0.133,0.545,0.133},
        stringstyle=\color[rgb]{0.627,0.126,0.941},
        captionpos=b,
        title=\lstname,
        showstringspaces=false
       }


\begin{document}

\begin{frame}
\titlepage
\end{frame}
 
 
%%%%%%%%%%%%%%%%%%%%%%%%%%%%%%%%%%%%%%%%%%%%%%%%%%%
\begin{frame}{Primary Objective}

Students will be able to 

\vfill

\begin{quote}
\textcolor{wcprimary}{define and use statically- and dynamically-sized arrays of numerical types in both C++ and Modern Fortran.
}
\end{quote}

\vfill 

\end{frame}


%%%%%%%%%%%%%%%%%%%%%%%%%%%%%%%%%%%%%%%%%%%%%%%%%%%
\begin{frame}{The Plan}

{\small
\begin{tabular}{ll}
 (01) C++ \& Fortran Toolchain       & (07) Shared-Memory Parallelism on CPUs (OpenMP) \\
 (02) C++ \& Fortran Basics          & (08) Shared-Memory Parallelism on GPUs (CUDA) \\
 (03) C++ \& Fortran Arrays          & (09) Distributed-Memory Parallelism (MPI, locally) \\
 (04) Matrix-Vector                  & (10) Distributed-Memory Parallelism (MPI, Beocat) \\
 (05) BLAS+LAPACK and FLOPs          & (11) Hybrid Parallelism (OpenMP+MPI, Beocat) \ \\
 (06) Overview of Parallel Computing & (12) Wrapping Up \& Future Trends \\
\end{tabular}
}
\vfill 


For today, execute \, {\tt sudo apt install valgrind}.

\end{frame}


%==============================================================================%
\section{C++ Arrays}
%==============================================================================%

%------------------------------------------------------------------------------%
\begin{frame}[fragile]{C++ Assignment}
 \lstinputlisting[language=C++]{compiled_code/assignment.cc}
\end{frame}

%------------------------------------------------------------------------------%
\begin{frame}[allowframebreaks]{C++ Arrays - Main Program}
 \lstinputlisting[language=C++]{compiled_code/basic_arrays.cc}
\end{frame}

%------------------------------------------------------------------------------%
\begin{frame}[allowframebreaks]{C++ Arrays - Function Header}
 \lstinputlisting[language=C++]{compiled_code/basic_arrays_functions.hh}
\end{frame}

%------------------------------------------------------------------------------%
\begin{frame}[allowframebreaks]{C++ Arrays - Function Definitions}
 \lstinputlisting[language=C++]{compiled_code/basic_arrays_functions.cc}
\end{frame}

%------------------------------------------------------------------------------%
\begin{frame}{Compiling}

 Go ahead, try {\tt g++ basic\_array.cc}.

 \pause
 
 \vfill
 The error is a {\bf linking error}.  
 
 \pause 
 
 \vfill
 Comment out the {\tt \#include "basic\_arrays\_functions.hh"} line 
 in {\tt basic\_array.cc} and try compiling again.
 
 \pause 
 \vfill
 The error is a {\bf syntactical error}---all names, including 
 functions, need to be defined before use.
 
 \pause 
 \vfill
 The right way (after uncommenting the {\tt .hh} file): \\
 
 {\tt g++ basic\_array\_functions.cc basic\_array.cc}\\
 
 (where the order matters!)
 
 
\end{frame}

%------------------------------------------------------------------------------%
\begin{frame}{Compiling with {\tt make}}
  \lstinputlisting[language=make]{compiled_code/make_basic_arrays}
\end{frame}

%------------------------------------------------------------------------------%
\begin{frame}[allowframebreaks]{Fortran Arrays - Main Program}
 \lstinputlisting[language=Fortran]{compiled_code/basic_arrays.f90}
\end{frame}

%------------------------------------------------------------------------------%
\begin{frame}[allowframebreaks]{Fortran Arrays - Functions Module}
 \lstinputlisting[language=Fortran]{compiled_code/basic_arrays_module.f90}
\end{frame}

%------------------------------------------------------------------------------%
\begin{frame}{Exercises}

Homework 8


\end{frame}


\end{document}

